\documentclass{article}
\usepackage{amsmath}
\usepackage{amssymb}
\usepackage{enumerate}

\begin{document}

\title{Physics 340 Final Exam: Mathematical Physics}
\author{}
\date{}
\maketitle

\section*{Instructions}
Answer all five questions. Show all work for full credit. Each question is worth 20 points.

\section*{Question 1: Algebra of Vectors (20 points)}

Consider the following two vectors in \( \mathbb{R}^3 \):
\[
\mathbf{A} = 2\hat{i} - \hat{j} + 4\hat{k}, 
\qquad
\mathbf{B} = -\hat{i} + 3\hat{j} + 2\hat{k}.
\]

\begin{enumerate}[a)]
\item Compute the dot product \( \mathbf{A}\cdot\mathbf{B} \).

\item Compute the cross product \( \mathbf{A}\times\mathbf{B} \).

\item Find the angle between the two vectors \( \mathbf{A} \) and \( \mathbf{B} \).
\end{enumerate}

\newpage
Continue your solution here
\vspace{4.0in}

\newpage

\section*{Question 2: Coupled Linear First-Order Differential Equations (20 points)}

Consider the coupled system:
\[
\frac{dx}{dt}=4x+y,
\qquad
\frac{dy}{dt}=2x+3y.
\]

\begin{enumerate}[a)]
\item Write the system in matrix form:
\[
\frac{d\vec{X}}{dt}=A\vec{X}.
\]

\item Find the eigenvalues and eigenvectors of the matrix \(A\).

\item Solve the system for \(x(t)\) and \(y(t)\), given
\[
x(0)=1,\qquad y(0)=2.
\]

\item Describe, in words, the long-term behavior of the system.
\end{enumerate}

\newpage
Continue your solution here
\vspace{4.0in}

\newpage

\newpage
Continue your solution here
\vspace{4.0in}

\newpage

\section*{Question 3: AC RC Circuit and Complex Impedance (20 points)}

An alternating voltage source
\[
V(t)=V_0 e^{i\omega t}
\]
is applied across a resistor \(R\) and a capacitor \(C\) connected in series.

\begin{enumerate}[a)]
\item Using complex impedance, derive the total impedance \(Z\) of the RC circuit in terms of \(R\), \(C\), and \(\omega\).

\item Using the complex form of Ohm's law,
\[
V(t)=I(t)Z,
\]
find an expression for the current \(I(t)\).

\item Find the phase shift between the voltage and current. State whether the current leads or lags the voltage.

\item Briefly explain the physical role of the resistor and the capacitor in producing this phase shift.
\end{enumerate}

\newpage
Continue your solution here
\vspace{4.0in}

\newpage

\newpage
Continue your solution here
\vspace{4.0in}

\newpage

\section*{Question 4: Fourier Analysis (20 points)}

Let \( f(x) \) be periodic with period \(2L\), defined on one period by
\[
f(x)=
\begin{cases}
-1, & -L<x<0,\\[6pt]
1, & 0<x<L.
\end{cases}
\]

\begin{enumerate}[a)]
\item Compute the Fourier coefficients \(a_0\), \(a_n\), and \(b_n\).

\item Write the Fourier series representation of \(f(x)\).

\item Discuss the convergence of the series at:
\begin{enumerate}[i)]
\item points where \(f(x)\) is continuous,
\item points where the periodic extension is discontinuous.
\end{enumerate}
\end{enumerate}

\newpage

\section*{Question 5: Vector Calculus (20 points)}

Let
\[
\mathbf{F}=(xy)\hat{i}+(x^2+z)\hat{j}+(y+z^2)\hat{k}.
\]

\begin{enumerate}[a)]
\item Compute the divergence:
\[
\nabla\cdot\mathbf{F}.
\]

\item Compute the curl:
\[
\nabla\times\mathbf{F}.
\]

\item Evaluate the line integral
\[
\int_C \mathbf{F}\cdot d\mathbf{r}
\]
along the curve
\[
\mathbf{r}(t)=\langle t,\; t^2,\; t\rangle,
\qquad 0\le t\le 1.
\]
\end{enumerate}

\newpage
Continue your solution here
\vspace{4.0in}

\newpage

\newpage
Continue your solution here
\vspace{4.0in}

\newpage

\end{document}
